%!TEX root = ../hutbthesis_main.tex

\chapter{基于运动模仿的生物合理肌肉骨骼运动控制}

\section{系统架构}\label{ux7cfbux7edfux67b6ux6784}

\subsection{整体架构}\label{ux6574ux4f53ux67b6ux6784}

Kinesis系统采用模块化的设计思想来实现运动模仿下的生物合理的肌肉骨骼运动控制，它的主要目的就是保证运动模仿的准确性以及生物力学的合理性。系统整体架构由数据处理、运动学、控制、仿真、评估五个主要部分组成，各个模块的职责清楚，接口统一，具体的模块划分和对应的实现文件见表~\ref{tab:3-1}。

\begin{table}[htbp]
\centering
\caption{Kinesis 系统核心模块说明}
\label{tab:3-1}
\small
\begin{tabularx}{\textwidth}{lX}
\toprule
模块 & 主要职责 \\
\midrule
数据处理模块 & 运动数据的加载、清洗、标准化、对齐、增强等预处理操作 \\
运动学模块 & 基于 SMPL 模型完成前向运动学计算，将关节角度转换为 3D 空间姿态，输出关节位置、速度等状态信息 \\
控制模块 & 构建 actor-critic 强化学习架构，基于 PPO 算法生成关节控制信号，优化运动控制策略 \\
仿真模块 & 基于 MuJoCo 物理引擎搭建肌肉骨骼系统仿真环境，执行控制信号并更新物理状态 \\
评估模块 & 从轨迹相似度、生物合理性、任务完成度等维度评估运动控制效果，反馈优化方向 \\
\bottomrule
\end{tabularx}
\end{table}

系统各个模块之间相互作用遵循的是“数据驱动、计算决策、仿真验证、评价反馈”的闭环逻辑，整个交互过程如图~\ref{fig:3-1}所示。数据处理模块输出的标准化运动数据给运动学计算提供输入，运动学模块把关节角度转成3D姿态，给控制模块提供状态感知的依据，控制模块产生控制信号，经过运动学转换之后输入到仿真模块中，仿真模块执行物理计算，得到环境状态，由评估模块计算奖励和损失，最后反馈给控制模块进行策略优化。

\begin{figure}[htbp]
\centering
\includegraphics[width=0.70\textwidth]{images/image3.png}
\caption{Kinesis 系统整体架构}
\label{fig:3-1}
\end{figure}

\subsection{模块间交互}\label{ux6a21ux5757ux95f4ux4ea4ux4e92}

数据处理流程，数据处理模块利用load\_motions方法读取KIT数据集，使用Python的multiprocessing库进行多线程并行预处理，完成数据清洗、标准化、对齐和增强之后，以结构化的字典形式保存关节角度、位置、速度等特征，供后面模块调用。

运动学计算交互，运动学模块接收到控制模块给出的关节角度指令之后，调用forward\_kinematics\_batch方法对关节进行前向运动学计算，得到关节的位置、速度、角速度等状态信息，给控制模块的策略决策提供实时的状态反馈。

控制决策交互部分，控制模块根据运动学模块输出的状态信息和目标运动轨迹，用Actor生成关节控制信号，经过运动学转换后得到符合
MuJoCo 关节顺序的角度指令，传送到仿真模块。

物理仿真交互是仿真模块在MuJoCo引擎中给定关节角度指令之后执行多关节的动力学计算和碰撞检测，改变肌肉骨骼系统物理状态，然后把当前的状态反馈给评价模块。

评估反馈交互就是根据预设的目标函数来计算轨迹误差、速度误差、生物力学约束惩罚和任务奖励，把评估结果反馈给控制模块，使PPO算法对策略网络参数进行优化。

模块化架构设计使得系统具有较好的扩展性，可以添加新的模块或者修改接口来适应不同的运动控制任务，而且模块之间耦合度低，利于系统维护。

\section{运动数据处理}\label{ux8fd0ux52a8ux6570ux636eux5904ux7406}

运动数据处理是实现高精度运动模仿的基础，它的好坏直接关系到策略学习的效果和评价的可靠性。系统由数据加载、预处理以及增强过程组成。

（1）运动数据加载与解析

系统采用标准的运动捕捉数据集作为输入，日常运动种类的数据被加载进去，数据加载的过程如下所示：

读取运动数据文件，提取关节旋转、根节点平移、帧率、运动时长等主要信息；

统一姿态表示方式，把不同的来源、不同的格式数据转化为系统可以兼容的结构；

按照动作类别对训练集和测试集进行划分，保证实验可以复现、比较。

加载完毕之后，系统就得到了结构化时序数据，包含了关节旋转、关节三维位置、根节点轨迹、线速度和角速度等全部的信息。

（2）运动数据预处理

原始运动捕获数据中存在着噪声、漂移、高度偏移、朝向不一致、帧率不统一等状况，只有经过标准化预处理之后才能被用作训练和评价的依据。具体的步骤有：

噪声滤波和平滑处理，用滤波算法消除高频抖动和异常跳变，使姿态序列连续平滑；

根节点高度校正，自动调节整体垂直位移，保证角色足部与地面良好接触，防止悬空或者穿模；

统一帧率、时间对齐，将所有运动序列重采样到相同的时间分辨率上，使时序对齐一致；

朝向归一化处理，将所有的动作初始运动方向统一，消除由于朝向不同造成的评估偏差；

异常姿态剔除就是检测出超出生理活动范围的姿态帧，防止仿真崩溃和异常行为的发生。

运动数据预处理流程如图~\ref{fig:3-2}所示。

\begin{figure}[htbp]
\centering
\includegraphics[width=0.92\textwidth]{images/image4.png}
\caption{运动数据预处理流程图}
\label{fig:3-2}
\end{figure}

（3）运动数据增强

为了提高策略的泛化能力和鲁棒性，防止过拟合，在训练过程中使用了许多数据增强技术：

随机朝向旋转，在水平面内随机旋转整体运动方向，保持动作内部相对关系不变；

随机噪声注入，给关节角度加上小幅度高斯噪声，模拟真实的运动自然波动；

时序偏移和速度缩放，随机截取一段运动片段或者小幅改变运动速度，提高策略对不同开始时间、运动速度的适应性；

局部姿态扰动就是给重要的关节施加很小的扰动，提高策略在非理想条件下稳定性。

数据增强可以扩大部分样本量，削弱模型对某些姿态、方向、速度的过分依赖，在没有见过的情况下也可以有较好的运动效果。

\section{运动模仿算法设计}\label{ux8fd0ux52a8ux6a21ux4effux7b97ux6cd5ux8bbeux8ba1}

（一）总体算法框架

算法采用经典 Actor-Critic 架构，结合 PPO 近端策略优化算法，整体框架如图~\ref{fig:3-3}所示。

\begin{figure}[htbp]
\centering
\includegraphics[width=0.92\textwidth]{images/image5.png}
\caption{运动模仿算法框架}
\label{fig:3-3}
\end{figure}

Actor以系统当前的状态作为输入，输出连续的、平滑的、有界的控制信号，直接驱动肌肉骨骼模型做动作。

Critic评价目前的状态所带来长期的收益，用来稳定策略的更新，减少梯度的波动，加快收敛。

奖励和损失函数是把轨迹误差、速度误差、生物合理性惩罚和任务完成奖励结合起来，使策略朝着更像、更稳、更合理的方向发展。

该架构对于高维、强耦合、连续动作的肌肉骨骼系统来说比较稳定，可以用来做复杂的运动模仿任务。

（二）强化学习算法改进（核心）

（1）改进型近端策略优化（PPO）

用裁剪目标函数来控制策略参数的变化幅度，防止策略参数突然发生大的变化，从而提高训练的稳定性以及样本的利用效率。

核心公式：
\begin{equation}
L^{CLIP}(\theta)=\min\left(r_i(\theta)A_i,\operatorname{clip}(r_i(\theta),1-\varepsilon,1+\varepsilon)A_i\right)
\tag{3-1}
\end{equation}

\(L^{CLIP}(\theta)\)：PPO
算法的裁剪损失函数。\(\theta\)是神经网络参数；这个损失用来更新策略，防止一步更新太大导致训练崩溃。

\(r_{i}(\theta)\)：旧的和新的策略的概率比。新的策略选当前动作的概率是以前策略选当前动作概率的商，即策略的变化幅度。

\(A_{i}\)：优势函数。表示当前动作比平均水平好多少，\(A_{i}\)\textgreater0动作好，要多强化；\(A_{i}\)\textless0动作差，要少用。

\(clip\left( r_{i}(\theta),1 - \varepsilon,1 + \varepsilon \right)\)：裁剪操作。把概率比rt限制在区间[1-$\epsilon$,~1+$\epsilon$]内，避免策略变化过猛。

\(\varepsilon\)就是裁剪系数，一般取0.2来控制策略更新幅度的上限。

\(\min\)：取两个值里较小的那个，保证更新保守、训练稳定。

（2）广义优势估计（GAE）优化

采用加权组合多步时序差分误差方法，即估计的精度与方差之间求平衡点，得到高精确的优势函数，策略更新更稳定。仅依靠
GAE
优化优势仍难以完全解决肌肉骨骼系统训练震荡、易发散的问题。本文进一步从梯度更新、经验复用完善训练流程。

在梯度更新方面，采用梯度裁剪 +
权重衰减组合策略，限制梯度最大范数，避免单次更新梯度爆炸；同时加入小权重衰减项，抑制网络参数过度增长，防止过拟合，让训练过程更平稳。

在经验复用机制上，采用优先级经验回放，优先采样高价值、高误差样本，提高样本利用率，减少无效样本干扰，加快收敛速度。

通过上述多维度训练优化，有效缓解肌肉骨骼动力学环境下训练震荡、策略发散、收敛慢等痛点，让算法在复杂运动模仿任务中更稳定、高效。

核心公式：
\begin{equation}
\widehat{A}_{t}^{GAE(\gamma,\lambda)}=\sum_{l=0}^{\infty}(\gamma\lambda)^l\delta_{t+l}
\tag{3-2}
\end{equation}

\({\widehat{A}}_{t}^{GAE}\)：广义优势函数，更稳、更准地衡量动作好坏。

$t$：当前时间步。

$\gamma$：折扣因子（未来奖励打折扣，常用 0.99）。

$\lambda$：GAE 平滑系数（权衡偏差与方差，常用 0.95）。

\((\gamma\lambda)^{l}\)：越往后权重越小，兼顾短期与长期。

\(\delta_{t + l}\)：第\(t + l\)步时序差分误差（实际奖励+后续价值-当前估计价值）。

含义是将未来若干步的好坏进行加权求和，优势更加稳定，训练不容易出现震荡。

（3）混合专家策略（MoE）

采用多专家网络、门控机制，根据不同的运动模式动态地把网络分配给不同的任务上，使复杂的动作得到更好的表达和泛化。

（三）策略网络优化

（1）高斯策略网络用高斯分布建模连续动作输出，学习均值和标准差来实现精准的肌肉激活控制，公式如下所示：
\begin{equation}
\pi_{\theta}(a\mid s)=\mathcal{N}\left(a;\mu_{\theta}(s),\sigma^2\right)
\tag{3-3}
\end{equation}

\(\pi_{\theta}(a \mid s)\)：策略，在状态\(s\)下输出动作\(a\)的概率分布。

\(\mathcal{N}\left( a;\mu_{\theta}(s),\sigma^{2} \right)\)：正态分布（高斯分布）。

\(a\)：动作（肌肉激活信号、关节指令）。

\(s\)：当前状态（关节角度、速度、参考轨迹等）。

\(\mu_{\theta}(s)\)：网络预测的动作均值（主输出）。

\(\sigma^{2}\)：方差（控制探索 / 随机程度）。

\(\theta\)：网络可训练参数。

含义为动作是带有均值和方差的连续随机变量，可以精确控制也可以随机探索。

（2）价值、策略网络都是用底层特征提取层来减少参数重复使用的情况、加快训练速度、收敛快慢，这两个方面。

（3）输出缩放和限幅就是把网络输出限定在生理安全范围内，防止出现指令突变、超调等情况，保证运动的平滑。

（四）目标函数设计

目标函数是算法优化的关键部分，既考虑模仿的准确性又考虑生物的合理性，采用多因素加权组合的形式，即模仿度权重×模仿误差+生物合理性权重×生物合理性。

\begin{equation}
L=w_1L_{\mathrm{pose}}+w_2L_{\mathrm{vel}}+w_3L_{\mathrm{bio}}+w_4R_{\mathrm{task}}
\tag{3-4}
\end{equation}
各项含义如下：

$L$：总损失（总代价），越小越好。

\(w_{1}\)为姿态损失权重，\(L_{\mathrm{pose}}\)为姿态或者关节位置误差损失。

\(w_{2}\)是速度损失权重，\(L_{\mathrm{vel}}\)是速度误差损失（快慢节奏不像就变大）。

\(w_{3}\)是生物约束权重，\(L_{\mathrm{bio}}\)是生物合理性惩罚，即关节超限、肌肉异常、耗能大就增大。

\(R_{\mathrm{task}}\)为任务奖励，即完成动作、保持平衡、不倒就变大，所以前面是减号。

含义是综合考虑像不像、稳不稳、合不合理、能不能成等几个方面来优化。

（五）策略网络设计

Actor是系统动作决策的核心，用轻量化、高泛化性、高运行稳定性的多层感知器（MLP）结构，在高维肌肉骨骼状态空间里可以快速输出平滑、连续、符合生理约束的控制信号。

输入层是系统目前所有可以观测的状态，即实时关节角度、关节运动速度、根节点位置和移动速度、目标参考轨迹等，给策略提供完整的运动感知依据。

隐藏层用全连接的多层非线性激活函数来完成高维时序特征的提取，在多层特征抽象的过程中慢慢捕捉运动时序的依赖关系、关节之间的耦合关系以及动态变化的规律，使得策略可以理解并复现复杂的运动模式。

输出层用有界激活函数把网络输出严格限定在安全数值范围内，再经过线性映射转换成肌肉骨骼模型可以直接执行的控制范围，保证控制信号连续、无突变、无超出生理限制的指令，保证运动过程平滑稳定、可以安全执行。

Actor所输出的控制信号会直接对肌肉骨骼驱动模型产生影响，从信号角度上保证运动姿态自然、发力平稳、过渡流畅，从而达到高精度、高可靠的模仿动作的目的。

（六）价值网络设计

Critic和Actor都使用共享特征提取层的模式，二者状态表示相同但是参数数量大大减少，训练效率提高，收敛速度加快。

输入信息同Actor输入完全一致，保证状态观测的一致性、无偏差，给价值评估提供准确的依据。

输出结构，输出单一标量价值，用来表征当前状态在未来时段内长期的期望收益，直接体现当前运动状态的好坏和风险程度。

训练目标就是最小化价值预测误差，用广义优势估计（GAE）来计算优势函数更加准确，减小策略更新时梯度的波动，提高整个训练的稳定性。

肌肉骨骼等高维、强耦合、强非线性复杂的系统中，Critic可以有效地抑制训练震荡、策略发散、仿真崩溃等问题，是稳定、可靠、可复现训练的重要组成部分。

（七）模仿损失计算

模仿损失是评价算法学习效果的主要指标，它直接体现生成的运动和示范的运动之间的差别，是由多个精细化的损失项加权得到的，共同促使策略朝着高质量模仿的方向发展。

位置损失主要是计算根节点、足部、手部等重要关节的三维空间位置误差，对平衡和姿态起着重要作用的部位给予更高的权重，保证空间轨迹高度符合。

姿态损失用四元数来计算关节的旋转差，防止欧拉角出现万向锁的问题，从全局姿态的角度保证动作整体结构的一致性。

速度损失就是对齐关节线速度和角速度，使生成的运动快慢节奏、动态变化和示范动作完全一致，提高动作的真实感。

根节点损失就是对整体的位移、朝向、高度进行约束，使运动不能出现漂移、倾斜、失衡的现象，保证全局稳定。

采用多维损失联合约束的方法，从粗糙到精细、从全局到局部来学习示范动作的空间结构、时序规律、动态特性等，获得高质量、高流畅度、稳定性好的运动模仿效果。

\section{生物合理性约束}\label{ux751fux7269ux5408ux7406ux6027ux7ea6ux675f}

为了使系统生成的运动在轨迹上接近示范动作，在生理机制、力学特性、运动习惯上都符合真实的运动规律，本文从关节活动、肌肉驱动、能量消耗、运动平滑性四个方面建立硬约束限制和软惩罚引导相结合的生物合理性约束体系，使运动控制的结果更加自然、安全、可信。

（1）关节角度约束

人体各个关节都有其固定的生理活动范围，超出这个范围就会出现反关节、过度屈伸、扭曲等非生理姿态。因此，系统给髋、膝、踝、脊柱、肩等重要关节设置出符合解剖学的安全角度范围，并且用硬约束保证安全加软惩罚引导最优的双重机制。硬约束在仿真中强制执行，关节角度超过生理极限就停止仿真，防止出现异常姿态；软惩罚在关节接近角度极限的时候逐渐增大惩罚值，使策略主动把关节运动保持在舒适、安全、自然的活动范围内，从根本上保证运动的生物合规性。

（2）肌肉驱动约束

肌肉驱动约束严格按照肌肉的力学特性以及神经控制规律来执行，使肌肉的激活模式更接近人体真实的收缩逻辑。系统利用多种规则来控制肌肉的激活信号不能超过人体所能达到的最大出力水平，也不能让肌肉的激活速度突然变大或者变小，不能出现瞬间的剧烈发力等非生理现象，保证肌肉长度一直处在有效的张力范围内，符合长度
-
张力的关系，还会对不合理的肌群协同模式进行惩罚，防止出现无效发力、冗余用力或者肌肉紊乱激活等问题，从而使得运动更加高效、协调、自然。

（3）能量消耗约束

真实的人体运动会尽量用最少的能耗、最高的效率来完成任务，所以能量消耗约束可以促使策略学习趋向于更省力、节能、可持续的运动方式。系统对频繁的、大幅度的、多余耗能的运动进行惩罚，使肌肉骨骼模型趋向平稳稳定的动作模式。经由能量优化之后，所形成的运动更加贴近人的本能习惯，而且可以缩减虚拟角色因持续工作而产生的疲惫感，从而加强长时间连续运动的稳定性和流畅性。

（4）运动平滑性约束

运动平滑性是提高动作真实性的重要因素，生硬突变、剧烈抖动、姿态跳变都会使运动显得机械、不自然。因此系统对关节加速度和加加速度（jerk）进行了严格的限制，防止运动出现突然的、突兀的变化；在控制信号输出端加入了滤波平滑处理，使连续的动作过渡更加自然；从算法上禁止了姿态跳变、非生理抖动、不连续切换，使整个运动呈现出人体特有的连贯、柔顺、协调的动态特征。

生物合理性约束体系如图~\ref{fig:3-4}所示。

\begin{figure}[htbp]
\centering
\includegraphics[width=0.92\textwidth]{images/image6.png}
\caption{生物合理性约束体系图}
\label{fig:3-4}
\end{figure}

利用多层、多维约束组合，在保证轨迹模仿精度的基础上，使得运动完全符合人体生理结构和运动规律，从而达到更准确、更流畅、更安全、更真实的动作效果，具有更好的现实应用价值。

\section{算法训练流程}\label{ux7b97ux6cd5ux8badux7ec3ux6d41ux7a0b}

本系统用PPO近端策略优化算法做端到端强化学习训练，整个过程高度标准化、参数可以固化、结果可以复现，在复杂的肌肉骨骼动力学环境中能够达到稳定的收敛效果，并且可以输出鲁棒性好、模仿精度高的控制策略。算法训练流程如图~\ref{fig:3-5}所示。

\begin{figure}[htbp]
\centering
\includegraphics[width=0.92\textwidth]{images/image7.png}
\caption{算法训练流程图}
\label{fig:3-5}
\end{figure}

（一）训练准备

训练开始前要完成全部的初始化工作，为稳定训练打下基础。具体包含完成物理仿真环境搭建，设置重力、摩擦系数、肌肉参数等物理条件，固定随机种子以保证实验可以重复，指定CPU或者GPU计算设备以及并行运算参数，加载经过标准化预处理和增强后的运动数据集，初始化Actor和Critic的权重参数，统一设置学习率、批量大小、训练迭代次数、PPO裁剪系数等关键超参数，保证训练过程稳定可控。

（二）训练执行流程

训练在闭环交互中迭代执行，分为四个核心步骤：

（1）轨迹采样就是让机器人按照当前策略在仿真环境里滚动执行运动，从而得到连续的状态序列、动作指令、奖励信号以及完整的轨迹信息，进而形成训练经验数据；

（2）优势估计采用广义优势估计（GAE），计算出的动作长期价值更加准确，梯度估计误差更小，训练稳定性更好；

（3）策略更新，用PPO裁剪目标函数对Actor和Critic进行多次更新，严格控制策略更新幅度，防止训练震荡、发散或者崩溃；

（4）评估保存就是每隔固定迭代周期对策略的性能进行量化测试，得到MPJPE、覆盖率、成功率等指标，并保存精度最高、稳定性最好的最优模型。

（三）训练方法

模型训练启动阶段，对策略网络与价值网络的权重参数统一运用正交初始化方式完成赋值，有效规避网络初始输出出现数值饱和的问题。仿真环境运行初始状态统一设定为标准中立姿态，彻底摒除初始状态随机波动带来的影响，让每一轮训练都能从相同基准起点开展实验。

参数优化环节选用线性递减式学习率调控方案，训练前期设置偏大的学习数值，助力智能体快速完成环境探索与行为试错；随着训练进程不断推进，循序渐进下调学习率数值，以此对动作策略开展精细化迭代优化。批次样本规模也配合训练阶段动态灵活改动，训练前期采用小批次样本训练模式，提升探索行为的响应速度；训练后期切换为大批次样本更新模式，保障参数迭代平稳可靠，统筹兼顾训练运行效率与模型收敛稳定性。
训练流程中增设提前停止判定机制，持续追踪验证数据集内平均关节位置误差、有效帧覆盖占比两项核心评价指标。若指标数值在连续多轮迭代中不再产生正向提升，系统将自动中止后续训练流程，既能杜绝模型过度拟合现象发生，也能减少硬件算力资源不必要消耗。按照固定迭代频次定期存储训练模型文件，妥善应对程序异常中断等突发状况，规避训练成果无故丢失的情况。

整个训练周期内不间断采集损失变化趋势、奖励收益波动、关节位置偏差等各类曲线数据，依托可视化监测手段，第一时间察觉梯度异常激增、决策策略性能下滑、肢体动作失控失效等各类训练异常问题。根据实际运行状态灵活调整超参数配置，全方位把控训练走向，保障整个训练过程平稳有序，同时实现实验结果可追溯、可重复复现。

（四）训练终止条件

训练结束后，系统就会给出最后得到的控制策略，可以用于运动模仿评价、可视化展示、鲁棒性检验和之后的实验研究。为兼顾模型训练质量与运算处理效率，本系统预先设置多项训练终止判定规则。当迭代轮次抵达预设最大值时，训练进程将自动终止；若轨迹偏差、场景覆盖率、任务完成成功率等核心评估参数，历经多轮迭代后数值逐步趋于平稳，不再出现明显提升波动，即可判定算法完成收敛，训练流程随之停止。同时支持人工干预操作，训练过程中出现异常状况时，可手动叫停运行程序。一旦内存、显存等硬件计算资源占用触及临界阈值，系统会强制终止程序运行，规避资源过载引发的故障问题，避免产生额外损耗。

训练流程完结后，系统会输出最终优化成型的控制决策方案。该策略可应用于动作复刻效果评估、动态仿真可视化演示、模型抗干扰能力测试等环节，也能够为后续相关对照实验与算法深度探究提供数据与模型支撑。
